%!TEX root = ../../../../exa-ma-d7.1.tex
\VThreeApplicationDraft{app-feelpp-ktirio-geom}
{Terrain generation}
{V. Chabannes; C. Prud'homme}
{WP1}{demonstrator}
{Generate a city mesh including elevation, roads, water bodies and land markers.}
{Feel++, CGAL, multithreaded execution}{benchmark-in-progress}

\paragraph{Versioned implementation evidence.}
At inspected commit \href{https://github.com/numpex/apps-feelpp/commit/bf1febf41cd615c20928e1bb4ef70678533c0587}{\texttt{bf1febf}},
the application repository contains Grenoble, Tokyo and Perlin terrain JSON inputs and
plot configurations under \texttt{src/ktirio-geom/terrain}.  This proves an
executable configuration exists, not that benchmark results are available.

\subsection{Ktirio-Geom implementation developments in 2026}
\label{sec:app:ktirio-geom:implementation-2026}

The 2026 developments in the CGAL-based Ktirio-Geom implementation cover
roof generation, tiled terrain construction, geographic-overlay processing,
surface remeshing and volume tetrahedralization. The terrain demonstrator
below exercises part of these capabilities; its timings concern terrain
generation rather than the full library. The implementation details complement
the methodological contribution described under WP1 Task~1.1.

\paragraph{Roof generation.}
Ridged roofs, including hipped, gabled, Mansard, Gambrel and Saltbox roofs,
are constructed by applying CGAL's two-dimensional constrained Delaunay
triangulation to interior straight skeletons. Vertices are elevated according
to their propagation speeds and the requested roof height; hipped roofs use
uniform weights. Pyramidal, dome and onion roofs use the footprint barycentre
and a prescribed apex. Intermediate vertices are generated from the footprint
vertices at the requested resolution and connected by a fan triangulation.

\paragraph{Tiled terrain generation.}
Elevation height matrices and overlay features, including road multilines,
building footprints and water polygons, are read and preprocessed using
multiple threads. A structured grid carrying elevation data is constructed
for each tile. Neighbour assignment and elevation ghost-cell computation
require synchronization. Contour lines are then extracted independently for
each tile using marching triangles with a prescribed elevation interval.
The contour polylines are inserted into CGAL's two-dimensional constrained
Delaunay triangulation and simplified using a configurable threshold;
the supplied configuration retains approximately 10\% of contour vertices.
Overlay features are clipped to the tile domain, preprocessed and inserted
as constraints, with their elevations interpolated from the terrain data.
A flood-fill traversal tags features according to constraint identifiers,
and each triangulation is converted to Ktirio's mesh data structure.
Tile meshes are merged at a synchronization stage. Postprocessing removes
duplicate points and autorefines triangles adjacent to tile boundaries,
while tile interiors are already watertight and connected. Vertex elevations
are then assigned and geometric transformations, including dilation, applied
when required.

\paragraph{Geographic-overlay processing.}
Multiline features are buffered with an offset determined by feature type
and class, such as highways or paths. Polygon repair is applied to the
buffered geometries, followed by Boolean union to combine components of
the same feature. An optional two-dimensional alpha-wrapping procedure from
CGAL merges geometries belonging to different features, for example railway
tracks or separate highway lanes.

\paragraph{Surface remeshing and sizing fields.}
CGAL's isotropic remeshing is integrated with custom sizing fields. Input
surfaces are converted to polygon soups, using property maps to retain
patches and constraints. Edge markers select the constraints to preserve:
constrained edges may be split but not collapsed, and mesh markers are
retained after remeshing. Four interfaces accept, respectively, a uniform
target edge length; minimum and maximum target lengths with a tolerance
for adaptive remeshing; pairs of triangle-marker names and target lengths
to assign different characteristic sizes to marked regions; or a functor
returning a target edge length for a mesh triangle.

\paragraph{Volume tetrahedralization and remeshing.}
An input surface mesh is converted to a polygon soup and tetrahedralized
using CGAL's Conforming Constrained Delaunay Triangulation in three dimensions.
A breadth-first traversal removes cells outside the domain and assigns
distinct subdomain indices to enclosed volumes. Optional tetrahedral
isotropic remeshing uses a prescribed target edge length. The resulting
tetrahedra are converted to Ktirio's mesh data structure.

\paragraph{Benchmark contract.}
\texttt{BM-FEEL-KTIRIO-01} covers single-node thread scaling and time-to-solution.
Within this single-node scope, strong scalability is evaluated by maintaining a fixed global problem size (the 81-tile Tokyo grid) while scaling the worker thread count ($t \in \{1, 2, 4, 8, 16, 32, 64\}$) to measure the reduction in time-to-solution.
Weak scalability is excluded for this specific use case.
All benchmarks were executed on a single 128-core compute node of the Gaya supercomputer.
Testing is intentionally capped at 64 threads to match the physical boundaries of a single CPU socket and the available task-level parallelism.
Measurements reflect a single execution pass per thread count, with raw timing metrics exported directly as JSON files.


\subsection{Description of the benchmark}

The benchmark assesses the ability of Ktirio-Geom to generate simulation-ready city meshes from heterogeneous geospatial data.
The objective is to characterize the execution time of the terrain generation pipeline and the scalability of its multithreaded processing stages on a single compute node.
The workflow is designed to construct a watertight and conforming 3D surface triangulation of a city environment by merging raw elevation height fields with GIS overlay data (such as buildings, roads, and water bodies).

The resulting mesh represents the terrain surface, with elevation and GIS features incorporated into the geometric representation.
Feature tags are preserved to support subsequent visualization, simulation preprocessing, and geometric analysis.


\subsection{Benchmarking tools used}

This benchmark was performed on the gaya supercomputer, using feelpp.benchmaring v4.0.0.
Performance measurements were collected using internal timing instrumentation integrated into Ktirio-Geom.
Each stage of the pipeline reports elapsed wall-clock time, with additional per-thread statistics for multithreaded sections.

For tile-level parallel operations, the benchmark reports the longest executing thread, reflecting the critical path and providing a conservative estimate of scalability limits.


\subsection{Input/Output Dataset Description}


\subsubsection{Input Data:}

The considered use case for this benchmark is the generation of a city mesh for Tokyo.
The domain consists of a $9 \times 9$ tile grid at zoom level 16, around $\sim 20 \text{ km}^2$.

The experiments were conducted on a single Gaya compute node equipped with 128 physical CPU cores. The benchmark uses up to 64 worker threads and does not rely on hyperthreading.

The input data consists of:

\begin{itemize}
    \item Elevation grids representing terrain height fields
    \item Overlay features in GIS format, for buildings, roads, water bodies, landuse, etc.
    \item Vertical separation between contour lines is set to 10m.
    \item The simplification threshold influencing contour line simplification is set to 0.1. Meaning that around 10\% of points from each contour line will be conserved.
\end{itemize}

\subsubsection{Output Data:}

The output consists of a conforming terrain surface triangulation, with tagged features (water,roads,buildings,... ) and timing metrics for each major processing stage. These outputs are suitable for visualization, simulation preprocessing, and further geometric analysis.

\begin{figure}[H]
  \centering
  \begin{minipage}{0.48\textwidth}
    \centering
    \includegraphics[width=\textwidth]{feelpp/feelpp-benchmark-terrain_mesh_tokyo_full.png}
    \subcaption{Full-domain view of the generated Tokyo terrain mesh.}
  \end{minipage}\hfill
  \begin{minipage}{0.48\textwidth}
    \centering
    \includegraphics[width=\textwidth]{feelpp/feelpp-benchmark-terrain_mesh_tokyo_detail.png}
    \subcaption{Detailed view of the generated triangulation.}
  \end{minipage}
  \caption{Generated terrain mesh for the Tokyo test case.}
  \label{fig:terrain_mesh_output}
\end{figure}

\subsection{Results summary}

The benchmark results show that the scalability of the terrain generation pipeline is primarily determined by the \textit{createMesh} stage. Increasing the number of worker threads reduces the execution time of the parallel portions of this stage, while the remaining processing stages exhibit limited scalability.

The \textit{createMesh} stage consists of three main steps:

\begin{enumerate}
    \item Pre-processing: Synchronization of neighboring tiles and preparation of the data structures required for mesh generation. This step is performed sequentially.
    \item Tile mesh generation: Each tile is processed independently to generate its local mesh. This step is parallelized across worker threads using a thread pool. The local meshes are generated using CGAL's constrained Delaunay triangulation package.
    \item Post-processing: The local meshes are merged into a single global mesh, which is then autorefined to ensure watertightness. The mesh is subsequently elevated and transformed to its global coordinate system. Autorefinement is parallelized, while elevation and affine transformations are performed sequentially.
\end{enumerate}


The benchmark results show that the scalability of the terrain generation pipeline is primarily determined by the \textit{createMesh} stage.
As shown in \Cref{fig:terrain_overall_times}, increasing the number of worker threads significantly reduces the execution time of this stage, while the other stages of the pipeline remain largely unchanged.

Within \textit{createMesh}, the main source of scalability is the \textit{createTileMeshes} step.
As shown in \Cref{fig:createmesh_breakdown}, its execution time decreases from 117.18s with one worker to 4.57s with 32 workers.
This corresponds to a measured speedup of 25.63$\times$, close to the ideal 32$\times$ speedup.
The corresponding speedup is shown in \Cref{fig:speedup_createtilemeshes}.

Beyond 32 workers, the performance no longer improves.
With 64 workers, \textit{createTileMeshes} takes 5.01s, resulting in a lower speedup of 23.37$\times$.
At this point, \textit{postProcess} becomes the dominant component of \textit{createMesh}, accounting for approximately 66\% of its execution time at both 32 and 64 workers.
Increasing the number of workers therefore provides no further benefit for the complete \textit{createMesh} stage in a 64-physical-core environment.


\pgfplotstableread[col sep=comma]{\detokenize{chapters/applications/specs/data/app-feelpp-ktirio-geom/TerrainTimes.csv}}\dataTerrainGenTimes
\pgfplotstableread[col sep=comma]{\detokenize{chapters/applications/specs/data/app-feelpp-ktirio-geom/TerrainCreateMeshTimes.csv}}\dataCreateMeshTimes
\pgfplotstableread[col sep=comma]{\detokenize{chapters/applications/specs/data/app-feelpp-ktirio-geom/TerrainCreateMeshTimesRelative.csv}}\dataTerrainCreateMeshTimesRelative
\pgfplotstableread[col sep=comma]{\detokenize{chapters/applications/specs/data/app-feelpp-ktirio-geom/TerrainCreateTileMeshesSpeedup.csv}}\dataTerrainCreateTileMeshesSpeedup

% ---------------------------------------------------------
% FIGURE 1: Overall Pipeline Times
% ---------------------------------------------------------
\begin{figure}[H]
  \centering
  \begin{tikzpicture}
    \begin{axis}[
      width=0.9\textwidth, height=7cm,
      ymin=0,
      xlabel={Max Workers},
      ylabel={Execution time (s)},
      xtick=data,
      xticklabels from table={\dataTerrainGenTimes}{max_workers},
      xticklabel style={font=\small},
      yticklabel style={font=\small},
      ymajorgrids=true, grid style={dashed, gray!40},
      legend columns=4,
      legend style={at={(0.5,-0.2)}, anchor=north, font=\small, draw=none}
    ]
      \addplot+[thick, mark=*, color=blue] table [x expr=\coordindex, y=createMesh] {\dataTerrainGenTimes};
      \addlegendentry{createMesh}

      \addplot+[thick, mark=square*, color=red] table [x expr=\coordindex, y=exportMesh] {\dataTerrainGenTimes};
      \addlegendentry{exportMesh}

      \addplot+[thick, mark=diamond*, color=orange] table [x expr=\coordindex, y=updateForUse] {\dataTerrainGenTimes};
      \addlegendentry{updateForUse}
    \end{axis}
  \end{tikzpicture}
  \caption{Terrain generation scaling performance across all processing stages.}
  \label{fig:terrain_overall_times}
\end{figure}

% ---------------------------------------------------------
% FIGURE 2: Grouped createMesh Breakdown (Absolute & Relative)
% ---------------------------------------------------------
\begin{figure}[H]
  \centering
  % Left Plot: Absolute Times
  \begin{minipage}{0.48\textwidth}
    \begin{tikzpicture}
      \begin{axis}[
        width=\textwidth, height=6.5cm,
        ymin=0,
        xlabel={Max Workers},
        ylabel={Execution time (s)},
        xtick=data,
        xticklabels from table={\dataCreateMeshTimes}{max_workers},
        xticklabel style={font=\scriptsize},
        yticklabel style={font=\scriptsize},
        ymajorgrids=true, grid style={dashed, gray!40},
        legend style={at={(0.5,-0.25)}, anchor=north, font=\scriptsize, draw=none}
      ]
        \addplot+[thick, mark=*, color=black] table [x expr=\coordindex, y=Total] {\dataCreateMeshTimes};
        \addlegendentry{Total}

        \addplot+[thick, mark=square*, color=red!80!black] table [x expr=\coordindex, y=createTileMeshes] {\dataCreateMeshTimes};
        \addlegendentry{createTileMeshes}

        \addplot+[thick, mark=triangle*, color=green!60!black] table [x expr=\coordindex, y=postProcess] {\dataCreateMeshTimes};
        \addlegendentry{postProcess}

        \addplot+[thick, mark=diamond*, color=blue!80!black] table [x expr=\coordindex, y=preProcess] {\dataCreateMeshTimes};
        \addlegendentry{preProcess}
      \end{axis}
    \end{tikzpicture}
  \end{minipage}\hfill
  % Right Plot: Relative Stacked Times
  \begin{minipage}{0.48\textwidth}
    \begin{tikzpicture}
      \begin{axis}[
        width=\textwidth, height=6.5cm,
        ymin=0, ymax=100,
        xlabel={Max Workers},
        ylabel={Relative time (\%)},
        xtick=data,
        xticklabels from table={\dataTerrainCreateMeshTimesRelative}{max_workers},
        xticklabel style={font=\scriptsize},
        yticklabel style={font=\scriptsize},
        ymajorgrids=true, grid style={dashed, gray!40},
        ybar stacked, bar width=12pt,
        legend style={at={(0.5,-0.25)}, anchor=north, font=\scriptsize, draw=none}
      ]
        \addplot+[ybar, fill=red!70, draw=black!70] table [x expr=\coordindex, y=createTileMeshes] {\dataTerrainCreateMeshTimesRelative};
        \addlegendentry{createTileMeshes}

        \addplot+[ybar, fill=green!70, draw=black!70] table [x expr=\coordindex, y=postProcess] {\dataTerrainCreateMeshTimesRelative};
        \addlegendentry{postProcess}

        \addplot+[ybar, fill=blue!70, draw=black!70] table [x expr=\coordindex, y=preProcess] {\dataTerrainCreateMeshTimesRelative};
        \addlegendentry{preProcess}
      \end{axis}
    \end{tikzpicture}
  \end{minipage}
  \caption{Breakdown of the \textit{createMesh} stage in absolute execution time (left) and relative proportions (right).}
  \label{fig:createmesh_breakdown}
\end{figure}

% ---------------------------------------------------------
% FIGURE 3: Speedup of createTileMeshes
% ---------------------------------------------------------
\begin{figure}[H]
  \centering
  \begin{tikzpicture}
    \begin{axis}[
      width=0.8\textwidth, height=6.5cm,
      ymin=0,
      xlabel={Max Workers},
      ylabel={Speedup},
      xtick=data,
      xticklabels from table={\dataTerrainCreateTileMeshesSpeedup}{max_workers},
      xticklabel style={font=\small},
      yticklabel style={font=\small},
      ymajorgrids=true, grid style={dashed, gray!40},
      legend style={at={(0.05,0.95)}, anchor=north west, font=\small}
    ]
      % Plot optimal and half-optimal bounds (No fillbetween library needed)
      \addplot[thick, dashed, color=black] table [x expr=\coordindex, y=optimal] {\dataTerrainCreateTileMeshesSpeedup};
      \addlegendentry{Optimal}

      \addplot[thick, dotted, color=black] table [x expr=\coordindex, y=half-optimal] {\dataTerrainCreateTileMeshesSpeedup};
      \addlegendentry{Half-optimal}

      % Plot actual speedup
      \addplot+[thick, mark=*, mark size=2pt, color=red!80!black] table [x expr=\coordindex, y=createTileMeshes] {\dataTerrainCreateTileMeshesSpeedup};
      \addlegendentry{createTileMeshes}

    \end{axis}
  \end{tikzpicture}
  \caption{Parallel speedup for the \textit{createTileMeshes} computation compared to theoretical boundaries.}
  \label{fig:speedup_createtilemeshes}
\end{figure}
